\documentclass[12pt,a4paper]{article}
\usepackage[utf8]{vntex}
\usepackage[utf8]{inputenc}
\usepackage{amsmath, amssymb, anyfontsize, array, color,  enumerate, enumitem, fancyhdr, float, graphicx, hyperref, lastpage, listings, mathtools, mips, natbib, pifont, tikz}
\usepackage[open,openlevel=0]{bookmark}
\usepackage[framemethod=tikz]{mdframed}

\usetikzlibrary{calc}
\newcommand\HRule{\rule{12cm}{1pt}}

\pagestyle{fancy}
\fancyhead{} % clear all header fields
\fancyhead[L]{
	\color{blue}
	\begin{tabular}{rl}
		\begin{picture}(25,15)(0,0)
		\put(0,-8){\includegraphics[width=8mm, height=8mm]{BK.png}}
		\end{picture}
		\begin{tabular}{l}
			\textbf{\bf \ttfamily Ho Chi Minh University of Technology}\\
			\textbf{\bf \ttfamily Faculty of Computer Science \& Engineering}
		\end{tabular}
	\end{tabular}
}
\fancyhead[R]{
	\begin{tabular}{l}
		\tiny \bf \\
		\tiny \bf
	\end{tabular}
}
\fancyfoot{} % clear all footer fields
\fancyfoot[L]{\scriptsize \ttfamily Computer Architecture Week 2}
\fancyfoot[R]{\scriptsize \ttfamily Page {\thepage}/\pageref{LastPage}}
\renewcommand{\headrulewidth}{0.3pt}
\renewcommand{\footrulewidth}{0.3pt}

\definecolor{dkgreen}{rgb}{0,0.6,0}
\definecolor{gray}{rgb}{0.5,0.5,0.5}
\definecolor{mauve}{rgb}{0.58,0,0.82}

\lstset{frame=tb,
	language=[MIPS]Assembler,
	aboveskip=3mm,
	belowskip=3mm,
	showstringspaces=false,
	columns=flexible,
	basicstyle={\small\ttfamily},
	numbers=none,
	numberstyle=\tiny\color{gray},
	keywordstyle=\color{blue},
	commentstyle=\color{dkgreen},
	stringstyle=\color{mauve},
	breaklines=true,
	breakatwhitespace=true,
	tabsize=3,
	numbers=left,
	stepnumber=1,
	numbersep=1pt,
	firstnumber=1,
	numberfirstline=true
}

\begin{document}
\begin{titlepage}
  \begin{tikzpicture}[remember picture, overlay]
    \draw[line width = 2pt] ($(current page.north west) + (1in,-1in)$) rectangle ($(current page.south east) + (-1in,1in)$);
  \end{tikzpicture}

  \begin{center}

    % Upper part of the page
    \textbf{\fontsize{12pt}{1pt}\selectfont HO CHI MINH CITY UNIVERSITY OF TECHNOLOGY}\\[0.5cm]
    {\fontsize{13pt}{1pt}\selectfont Faculty of Computer Science \& Engineering}\\[0.5cm]
    \begin{figure}[H]
      \centering
      \includegraphics[width=1.7in,height=1.7in]{BK.png}
    \end{figure}

    % Title
    \HRule\\[0.5cm]
    { \textbf{\fontsize{25pt}{1pt}\selectfont Computer Architecture Week 2}}\\[0.4cm]

    \HRule\\[0.8cm]
    \begin{minipage}{0.545\textwidth}
      \begin{flushleft}
        \textbf{Authors:}\\
        \begin{tabular}{l l}
          Lưu Nguyễn Hoàng Minh & 1952845 \\
          Vũ Anh Nhi            & 1952380 \\
          Nguyễn Chính Khôi     & 1952793 \\
          Nguyễn Hoàng          & 1952255 \\
        \end{tabular}
      \end{flushleft}
    \end{minipage}
    \begin{minipage}{0.4\textwidth}
      \begin{flushright}
        \textbf{Class:}\\
        CO2008\_CC02\\
        \textbf{Lecturer:}\\
        Kiều Đỗ Nguyên Bình\\

      \end{flushright}
    \end{minipage}

    \vfill

    % Bottom of the page
    \vspace{2cm}
    {\large} %{\large \today}
  \end{center}
\end{titlepage}


\pdfbookmark[0]{Question 1}{ques1}
\section*{Question 1:}
\begin{enumerate}[label=\alph*.]
  \item\mbox{}\\
        \begin{mdframed}[hidealllines=true,backgroundcolor=magenta!10]
          \begin{lstlisting}
  # A demonstration of some simple MIPS instructions
  # used to test QtSPIM

    # Declare main as a global function
    .globl main

    # All program code is placed after the
    # .text assembler directive
    .text

  # The label 'main' represents the starting point
  main:
    li $t2, 25		# Load immediate value (25)
    lw $t3, value		# Load the word stored in value (see bottom)
    add $t4, $t2, $t3	# Add
    sub $t5, $t2, $t3	# Subtract
    sw $t5, Z		#Store the answer in Z (declared at the bottom)

    # Exit the program by means of a syscall.
    # There are many syscalls - pick the desired one
    # by placing its code in $v0. The code for exit is "10"
    li $v0, 10 # Sets $v0 to "10" to select exit syscall
    syscall # Exit

    # All memory structures are placed after the
    # .data assembler directive
    .data

    # The .word assembler directive reserves space
    # in memory for a single 4-byte word (or multiple 4-byte words)
    # and assigns that memory location an initial value
    # (or a comma separated list of initial values)
  value:	.word 12
  Z:	.word 0
          \end{lstlisting}
        \end{mdframed}
\end{enumerate}

\pdfbookmark[0]{Question 2}{ques2}
\section*{Question 2:}

\pdfbookmark[0]{Question 3}{ques3}
\section*{Question 3:}

\pdfbookmark[0]{Question 4}{ques4}
\section*{Question 4:}

\pdfbookmark[0]{Question 5}{ques5}
\section*{Question 5:}

\pdfbookmark[0]{Question 6}{ques6}
\section*{Question 6:}

\end{document}
